\bookchapter{Chapitre 13 : Recommandations d'applications utilitaires}

\chapterintro{Ce chapitre sélectionne des outils de productivité, de santé et de voyage, et complète les capacités système et multi-appareils : nettoyage et anti-indésirables (Qingzhuo…), interrupteurs de blocage intégrés, météo et alertes sismiques ; contrôle à distance d'ordinateur (Sunlogin/ToDesk), bureautique IA (Trae…), bureaux animés, transfert local rapide (LocalSend), presse-papiers multi-appareils, codes dynamiques (Aegis/Microsoft Authenticator), planification TickTick et numérisation de documents papier.}

\booksection{I. Outils de productivité : gérer votre temps et votre attention}

Les outils de productivité sont l'« ossature » de la vie numérique : ils vous aident à mettre de l'ordre dans idées, tâches et documents épars. Recommandations par usage ci-dessous.

\booksubsection{Listes de tâches (gestion des tâches)}

\begin{itemize}
  \item TickTick (Dida Qingdan) — Pour qui : ceux qui veulent centraliser dans un même endroit leurs tâches et agendas de travail, de vie et d'études ; l'un des coups de cœur de cette section. Points forts : synchronisation entre téléphone/ordinateur/tablette/Web ; listes + calendrier + minuteur Pomodoro en un seul outil ; saisie en langage naturel — « réunion demain à 15 h », « rapport hebdo chaque lundi » — avec reconnaissance automatique des dates et des récurrences ; sous-tâches, étiquettes, priorités, pièces jointes et listes collaboratives. Idéal pour une gestion tout-en-un de « planning + suivi des tâches + revue quotidienne ».
  \item Bien utiliser TickTick : Boîte de collecte : tout jeter d'abord dans « Aujourd'hui / Collecte », puis trier le soir entre travail, famille et études.
  \item Vue calendrier : faites glisser réunions et échéances dans le calendrier pour voir les conflits, plus visuel qu'une simple liste.
  \item Tâches récurrentes : médicaments, factures, sport — définissez-les une fois, sans ressaisie.
  \item Pomodoro : lancez la concentration directement sur une tâche et voyez ce que vous avez réellement accompli dans la semaine.
  \item Liaison avec le téléphone : via l'assistant vocal « Rappelle-moi… », les rappels entrent dans le système ; pour les échéances importantes, doublez avec le calendrier système.
\end{itemize}

Todoist Public : utilisateurs minimalistes et multi-plateformes. Points forts : interface épurée, hiérarchie claire (tâche sous-tâche sous-sous-tâche), karma pour encourager les bonnes habitudes.

Microsoft To Do Public : utilisateurs professionnels et Outlook. Points forts : gratuit, sans publicité, fonction « Ma journée » qui se vide chaque jour pour alléger la charge mentale, intégration profonde à l'écosystème Microsoft.

Notes/Rappels système Public : utilisateurs légers qui ne veulent pas d'applications tierces. Points forts : expérience native, zéro latence, ajout via Siri/Xiao Ai, widgets système.

\booksubsection{Minuteur Pomodoro (entraînement de la concentration)}

\begin{itemize}
  \item Forest Public : ceux qui décrochent trop souvent de leur téléphone. Points forts : plantez un arbre en restant concentré ; si vous quittez l'application trop tôt, l'arbre meurt. Retour positif par la « vie virtuelle », plantation possible entre amis.
  \item Tide (Chaoxi) — Pour qui : ceux qui ont besoin de bruit blanc pour se concentrer ou s'endormir. Points forts : sons naturels (pluie, vagues, forêt) + Pomodoro + méditation réunis dans une seule appli, avec une interface épurée et apaisante.
  \item Principe : classiquement 25 minutes de concentration + 5 minutes de pause ; adaptez à votre état (50/10 ou 90/20).
\end{itemize}

\booksubsection{Calendrier (visualiser le temps)}

\begin{itemize}
  \item Calendrier système Public : tous les utilisateurs. Points forts : notifications natives stables, très faible consommation batterie, ajout via Siri/Xiao Ai.
  \item Google Calendar Public : multi-appareils et partage d'emploi du temps. Points forts : la fonction « Objectifs » planifie automatiquement sport/études ; Gmail détecte les événements dans les e-mails.
\end{itemize}

\booksubsection{Notes (capitaliser le savoir)}

\begin{itemize}
  \item Notion Public : unifier notes, tâches, bases de données et wiki ; étudiants, indépendants, équipes. Points forts : base de connaissances polyvalente, édition par blocs, pages personnelles, listes de lecture, tableaux de projets, riches modèles.
  \item Obsidian Public : gestionnaires de connaissances avancés, chercheurs, rédacteurs. Points forts : stockage Markdown local, liens bidirectionnels et graphe de relations — un « second cerveau » dont vous restez propriétaire.
  \item flomo Public : noter des idées brèves sans rédiger de longs textes. Points forts : notes en cartes, comme un statut, classification par \#étiquettes, consultation par chronologie.
  \item Notes Apple Public : utilisateurs iPhone/Mac avec des besoins légers. Points forts : lancement rapide, numérisation de documents, dessin, listes, synchronisation iCloud ; gratuit et parfaitement intégré.
  \item Evernote (Yinxiang Biji) — Pour qui : les utilisateurs avancés qui capturent beaucoup de pages web et archivent leurs documents. Points forts : une extension de capture web très puissante, de l'OCR pour lire le texte des images, un service vétéran et stable — mais la version premium est payante.
\end{itemize}

\booksubsection{Outils PDF}

\begin{itemize}
  \item PDF Expert Public : utilisateurs iPad/iPhone, professionnels annotant des contrats, étudiants. Points forts : reconnu comme l'un des meilleurs éditeurs PDF sur iOS, annotations fluides, signature, fusion/scission, formulaires.
  \item WPS Public : utilisateurs devant traiter Word/Excel/PPT/PDF. Points forts : suite bureautique tout-en-un, bonne compatibilité avec les stockages chinois et les fichiers WeChat ; la version gratuite suffit au quotidien.
  \item Adobe Acrobat Public : exigeants sur la norme PDF (droit, design, édition). Points forts : standard officiel, conversion PDF/A fiable, synchronisation multi-appareils et signature électronique professionnelles.
\end{itemize}

\booksubsection{Numérisation (papier numérique)}

Contrats papier, factures, devoirs, notes de tableau blanc, copies d'identité… l'appareil photo du téléphone remplace déjà la plupart des scanners. Processus : photographier nettement recadrage/redressement automatique export PDF/image OCR facultatif.

\begin{itemize}
  \item CamScanner Public : professionnels et étudiants scannant souvent contrats, reçus et pièces d'identité. Points forts : détection automatique des bords, correction de perspective, OCR, PDF multi-pages et synchronisation cloud.
  \item Microsoft Lens Public : étudiants, conversion de tableaux blancs ou de pages de livre en Word. Points forts : gratuit sans publicité, export en un clic vers OneNote/Word, reconnaissance précise.
  \item Numérisation dans Notes Apple Public : utilisateurs iPhone avec besoins légers. Points forts : aucune application à installer, icône appareil photo dans Notes, suppression automatique des ombres, signature, gratuit et efficace.
  \item Baimiao / CamScanner Web / mode « Document » des appareils photo de marque : les caméras système Huawei, Xiaomi, OPPO, vivo intègrent souvent un mode « Document/Scan » ; sur Android, Google Camera ou le scan système conviennent aussi. Parfait pour numériser en urgence une page de contrat.
  \item Conseils photo : lumière uniforme, téléphone parallèle au papier, cadrage des quatre coins ; pour plusieurs pages, photographiez dans l'ordre puis fusionnez en PDF ; pour les documents sensibles, pensez au filigrane et à l'emplacement de stockage, puis videz « Récemment supprimés ».
\end{itemize}

\booksection{II. Santé et sport}

Les applications de santé sont l'assistant personnel le plus sous-estimé du téléphone : elles suivent en continu sommeil, fréquence cardiaque et activité, bien mieux qu'une visite occasionnelle à l'hôpital pour refléter l'état de long terme.

\booksubsection{Sport et forme}

\begin{itemize}
  \item Keep Public : entraînement à domicile, débutants, perte de poids et mise en forme. Points forts : très nombreuses formations gratuites (HIIT, yoga…), plan d'entraînement généré par IA, communauté dynamique.
  \item Nike Training Club (NTC) Public : ceux qui veulent un coaching professionnel et structuré. Points forts : application officielle Nike, cours gratuits, démonstrations par des coachs, classement clair par muscle/équipement.
  \item Apple Fitness Public : utilisateurs Apple Watch motivés par les anneaux d'activité. Points forts : trois anneaux quotidiens (Bouger/Exercice/Se lever), objectifs visuels, données de fréquence cardiaque de la montre.
\end{itemize}

\booksubsection{Suivi du sommeil}

\begin{itemize}
  \item AutoSleep Public : utilisateurs Apple Watch voulant connaître leur sommeil profond/léger. Points forts : suivi automatique en dormant avec la montre, sans démarrage manuel, score de qualité et conseils de coucher.
  \item Sleep Cycle Public : ceux qui dorment sans montre et veulent surveiller avec le téléphone. Points forts : analyse des mouvements et du ronflement via le micro, réveil intelligent en sommeil léger pour un lever plus doux.
  \item Mi Fitness Public : utilisateurs de bracelets/montres Xiaomi. Points forts : gratuit et parfaitement intégré, suivi cœur/oxygène/sommeil, excellent rapport qualité-prix.
\end{itemize}

\booksubsection{Méditation et détente}

\begin{itemize}
  \item Calm Public : salariés stressés, difficultés d'endormissement. Points forts : méditation guidée « Daily Calm » + histoires du soir (voix célèbres), interface soignée, riche contenu anglais.
  \item Tide (Chaoxi) — Pour qui : les utilisateurs sinophones qui veulent du bruit blanc pour se concentrer ou s'endormir. Points forts : une perle chinoise qui réunit sons naturels + cours de méditation + Pomodoro, avec des guides en chinois plus chaleureux.
  \item Xiao Sleep — Pour qui : ceux qui veulent des sons pour s'endormir ou enregistrer leurs somniloquies et ronflements. Points forts : une immense banque de bruits blancs à superposer (pluie + ventilateur + feu de camp), un arrêt programmé au coucher, et l'enregistrement des propos nocturnes, à la fois drôle et utile.
\end{itemize}

\booksubsection{Santé féminine}

\begin{itemize}
  \item Flo Public : suivi des règles, projet de bébé, santé féminine. Points forts : prédictions IA précises des prochaines règles/ovulation, questions anonymes dans la communauté, design élégant.
  \item Meiyou — Pour qui : les utilisatrices en Chine et celles qui veulent des communautés et de la vulgarisation en chinois. Points forts : des contenus localisés très riches, avec des informations adaptées à chaque étape — cycles, projet bébé, grossesse, puériculture.
\end{itemize}

\booksubsection{Rappels de médicaments}

\begin{itemize}
  \item Yongyao Zhushou (Assistant Médicaments) — Pour qui : les malades chroniques sous traitement au long cours et les enfants qui gèrent les médicaments de leurs parents âgés. Points forts : identification des médicaments par scan du code-barres, liste de traitement créée automatiquement, rappels à l'heure par notifications, et journal de prise partageable avec la famille.
\end{itemize}

\booksubsection{Centre de données}

\begin{itemize}
  \item Apple Santé / Google Fit Public : tous les utilisateurs de smartphones. Points forts : agrège les données du téléphone, du bracelet et des applications tierces (pas, cœur, sommeil, règles…) en un « dossier de santé » unifié.
\end{itemize}

\booksection{III. Voyage et déplacements}

En voyage, quelques applications bien choisies évitent bien des ennuis. Recommandations par étape du déplacement.

\booksubsection{Avions et trains à grande vitesse}

\begin{itemize}
  \item Umetrip (Hanglv Zongheng) — Pour qui : les professionnels qui volent souvent. Points forts : données officielles de l'aviation civile — la source la plus fiable pour le suivi des vols, les changements de porte d'embarquement et les alertes de retard —, plus le type d'appareil et le plan de cabine consultables.
  \item 12306 Public : tous les voyageurs en train en Chine. Points forts : canal officiel d'achat de billets, liste d'attente automatique en cas d'annulation, billets étudiants/groupes pratiques.
  \item VariFlight (Feichangzhun) — Pour qui : les voyageurs sensibles à la ponctualité des vols et ceux qui déposent ou récupèrent quelqu'un à l'aéroport. Points forts : des prévisions de retard très justes, et le partage d'itinéraire pour que ta famille suive en direct ta position et ton heure d'arrivée.
\end{itemize}

\booksubsection{Cartes et navigation}

\begin{itemize}
  \item Amap (Gaode Ditu) — Pour qui : les citadins en Chine, conducteurs et habitués des VTC. Points forts : un trafic en temps réel très précis, une vraie capacité à éviter les bouchons, course ou covoiturage réservés au même endroit, et de nombreuses propositions de correspondances en transports en commun.
  \item Baidu Maps Public : besoin d'informations de lieux et de navigation intérieure. Points forts : vues panoramiques / navigation AR, plans détaillés de centres commerciaux et sites, nombreuses voix.
  \item Google Maps Public : voyageurs à l'étranger. Points forts : couverture mondiale, avis authentiques, Street View et cartes hors ligne indispensables hors de chez soi.
\end{itemize}

\booksubsection{Traduction}

\begin{itemize}
  \item Google Translate Public : voyageurs et apprenants multilingues. Points forts : traduction bidirectionnelle en mode « Conversation », traduction instantanée par l'appareil photo, packs hors ligne.
\end{itemize}

\booksubsection{Conversion de devises}

\begin{itemize}
  \item Jijian Huilü (convertisseur minimaliste) — Pour qui : les voyageurs à l'étranger et les adeptes des achats transfrontaliers. Points forts : une interface ultra épurée, plus de 160 devises, des taux actualisés en temps réel, la comparaison simultanée de plusieurs devises et une conversion rapide par appui long sur le chiffre.
\end{itemize}

\booksubsection{Météo}

\begin{itemize}
  \item Application Météo système / China Weather Public : tout le monde, prévisions et alertes officielles. Points forts : les applications météo constructeurs et canaux officiels comme « China Weather » relaient les alertes de l'observatoire — pluies torrentielles, typhons, fortes chaleurs — par bannière/notification, avec une forte autorité.
  \item Caiyun Weather Public : prévisions de pluie à la minute. Points forts : « va-t-il pleuvoir dans les deux heures ? » très précis, carte dynamique des précipitations, alerte avant la pluie.
  \item Carrot Weather Public : amateurs de météo aimant l'humour et les widgets riches. Points forts : bulletins météo par une IA mordante, nombreux widgets personnalisables, radar professionnel ; la version payante est plus complète.
  \item Points clés des alertes : autorisez les notifications des applications météo ; les alertes vont du bleu au rouge, le rouge indiquant le plus haut risque — adaptez vos déplacements. Voir aussi la section 7 de ce chapitre « Alertes météo / catastrophes naturelles / séismes ».
\end{itemize}

\booksubsection{Guides de voyage}

\begin{itemize}
  \item Mafengwo Public : amateurs de guides détaillés, road trips et voyages approfondis. Points forts : récits originaux riches, itinéraires/astuces/pièges à éviter complets, guides hors ligne téléchargeables.
  \item Xiaohongshu (RedNote) — Pour qui : les jeunes voyageurs en quête de retours d'expérience authentiques, de lieux hors des sentiers battus et de bonnes tables. Points forts : un format photos et courtes vidéos très parlant, des guides comme « 3 jours dans la ville X » vite mis à jour et toujours d'actualité, et une recherche très agréable.
\end{itemize}

\booksection{IV. Premiers pas avec l'automatisation}

L'automatisation est un levier « une configuration, un long confort » : elle peut réduire à zéro les dizaines de gestes répétés chaque jour. Cette section n'en donne que le concept ; pour les usages détaillés, voir le chapitre 23 : Pensée informatique et initiation à l'automatisation.

\begin{itemize}
  \item Raccourcis iOS : outil d'automatisation officiel d'Apple, sans programmation. En un geste, enchaînez par exemple « à la maison : activer le Wi-Fi + baisser la luminosité + annoncer la météo ».
  \item Automatisation Android : les scénarios système (« Tâches automatiques » de Xiaomi, « Intelligence de situation » de Huawei) sont les plus simples à prendre en main ; Tasker est très puissant mais exigeant, réservé aux geeks.
  \item Idée clé : repérez d'abord « les gestes répétés chaque jour », puis demandez-vous « un déclencheur + plusieurs actions peut-il les remplacer ? ». Exemples détaillés et usages avancés au chapitre 23.
\end{itemize}

\booksection{V. Navigateurs et moteurs de recherche : votre porte d'entrée sur l'information}

Le navigateur est votre fenêtre sur Internet ; le moteur de recherche, votre boussole dans l'océan d'informations. Bien les utiliser vous évite bien des détours.

\booksubsection{Comment fonctionne un navigateur :}

Lorsque vous saisissez une adresse, le navigateur : (1) interroge un serveur DNS pour obtenir l'adresse IP du nom de domaine ; (2) envoie une requête HTTP/HTTPS à ce serveur ; (3) le serveur renvoie le code de la page (HTML, CSS, JavaScript) ; (4) le navigateur interprète ce code et affiche la page.

\booksubsection{Navigateurs recommandés :}

\begin{itemize}
  \item Microsoft Edge : basé sur Chromium, rapide et compatible. Assistant Copilot intégré pour résumer, traduire et dialoguer. Synchronisation multi-appareils ; notre navigateur principal recommandé.
  \item Safari : navigateur par défaut sur Apple, profondément intégré, économe en énergie. Bonne protection de la vie privée, mais moins d'extensions.
  \item Chrome : le plus utilisé au monde, riche en extensions. Mais gourmand en mémoire et relativement énergivore sur mobile.
\end{itemize}

\booksubsection{Choix du moteur de recherche :}

\begin{itemize}
  \item Bing : moteur de Microsoft, étroitement intégré à Edge. Les résumés IA donnent directement la réponse, pas seulement des liens ; très efficace. Recherche d'images également correcte.
  \item Yandex : principal moteur russe, très fort en recherche d'images et recherche inversée. Idéal pour trouver des images similaires ou identifier un objet.
  \item Baidu : acteur majeur de la recherche en Chine, mais beaucoup de publicités ; un esprit critique est nécessaire.
  \item Google : le moteur le plus précis au monde, mais nécessite un accès réseau particulier en Chine.
\end{itemize}

\booksubsection{Guide de recherche d'information efficace :}

\begin{itemize}
  \item Utilisez la recherche IA : les résumés de Bing et le Copilot d'Edge répondent directement et extraient l'essentiel — bien plus rapide que d'ouvrir chaque lien.
  \item Repérer fausses informations et publicités : les résultats marqués « Publicité » ou « Sponsorisé » sont des promotions payantes, pas des résultats naturenels. Vérifiez l'autorité de la source et privilégiez sites officiels et médias reconnus.
  \item Télécharger correctement une application : tapez « nom de l'app + site officiel » (par exemple « WeChat site officiel ») et cliquez d'abord sur les résultats avec le domaine officiel. Évitez les sites de téléchargement tiers, souvent assortis de logiciels indésirables.
  \item Reconnaître un site officiel : pour Windows 11, allez sur microsoft.com ; pour Chrome, sur google.com. Vérifier que le domaine correspond à la marque est la méthode la plus élémentaire.
\end{itemize}

\booksubsection{Comprendre les noms de domaine :}

Le nom de domaine est l'« adresse » d'un site ; le lire aide à juger sa nature et son origine :

\begin{itemize}
  \item Extensions courantes : .com — organisation commerciale, la plus courante
  \item .cn — domaine national de la Chine
  \item .org — organisation à but non lucratif
  \item .net — services liés aux réseaux
  \item .edu — établissement d'enseignement
  \item .gov — organisme gouvernemental
  \item .xyz/.tech/.info — extensions récentes, faible barrière à l'entrée, à aborder avec prudence
\end{itemize}

Domaines nationaux/territoriaux : .jp — Japon

.kr — Corée

.uk — Royaume-Uni

.de — Allemagne

.tw — Taïwan

.hk — Hong Kong

Sécurité : les sites .gov et .edu sont en général plus fiables, car l'obtention est plus exigeante. Les domaines bon marché comme .xyz ou .top sont souvent détournés par des sites d'hameçonnage — restez vigilant.

\booksubsection{Blocage des publicités :}

\begin{itemize}
  \item AdGuard Public : tous les utilisateurs harcelés par pop-ups, bannières et publicités in-app. Points forts : blocage au niveau système (pas seulement le navigateur), filtrage des bannières et du pistage au démarrage des applications, économie de données et chargement plus rapide ; la version payante couvre tous les appareils.
\end{itemize}

\booksubsection{Navigateurs axés vie privée :}

\begin{itemize}
  \item Brave Public : ceux qui protègent leur vie privée et refusent le pistage publicitaire. Points forts : base Chromium, blocage natif publicités et traqueurs, mode Tor intégré, plus rapide que Chrome et moins gourmand en mémoire.
  \item DuckDuckGo Public : ceux qui ne veulent pas être « traqués par des pubs » après une simple recherche. Points forts : moteur éponyme (pas d'historique), effacement en un clic de toutes les traces (« Fire Button »), expérience cohérente sur toutes les plateformes.
\end{itemize}

\booksubsection{Tableau comparatif des navigateurs :}

\begin{center}
\begin{tabularx}{\linewidth}{|X|X|X|}
\hline
\textbf{Navigateur} & \textbf{Points forts} & \textbf{Public} \tabularnewline
\hline
Microsoft Edge & Noyau Chromium + Copilot IA intégré : résumé de pages, traduction, conversation & Professionnels axés efficacité, grands utilisateurs d'IA (recommandé en navigateur principal) \tabularnewline
Safari & Intégration profonde à l'écosystème Apple, économe, bonne confidentialité & Utilisateurs iPhone/Mac attentifs à l'autonomie et à l'intégration système \tabularnewline
Chrome & Le plus utilisé au monde, écosystème d'extensions le plus riche & Utilisateurs avancés ayant besoin de nombreuses extensions, synchronisation multi-appareils \tabularnewline
Brave & Blocage natif publicités et traqueurs, plus rapide & Utilisateurs soucieux de leur vie privée qui veulent bloquer les pubs sans extensions \tabularnewline
DuckDuckGo & Moteur éponyme sans historique, effacement des traces en un clic & Utilisateurs sensibles à la vie privée, opposés au pistage personnalisé \tabularnewline
\hline
\end{tabularx}
\end{center}

\begin{tcbpractice}{Essayez de créer une automatisation « détente du soir »}
\begin{enumerate}
  \item Créez le déclencheur : une heure précise (par exemple 22 h 30).
  \item Ajoutez les actions : baisser le volume > activer Ne pas déranger > lire une playlist détente > baisser la luminosité.
  \item Enregistrez et testez.
\end{enumerate}

\end{tcbpractice}

\booksection{VI. Outils système pratiques : nettoyage, anti-indésirables et sécurité}

Au-delà des applications de productivité et de santé, l'expérience quotidienne dépend surtout de trois choses : « être moins dérangé », « avoir assez d'espace de stockage » et « laisser les escrocs dehors ». Nous passons d'abord en revue les logiciels pratiques, puis les fonctions de blocage intégrées au téléphone — nombre de propriétaires ne les ont jamais activées.

\booksubsection{Nettoyage système et optimisation du stockage}

\begin{itemize}
  \item Qingzhuo — Pour qui : les utilisateurs Android qui veulent un nettoyage léger des caches, fichiers résiduels et dossiers vides, sans subir la batterie d'applis des « gestionnaires de téléphone ». Points forts : léger et sans publicité, avec scan approfondi des résidus d'applis, localisation des gros fichiers et analyse de l'occupation du stockage ; plus mesuré que les « accélérateurs en un clic » du commerce, parfait pour un entretien régulier si tu aimes mettre les mains dans le cambouis. Attention : vérifie les chemins avant de nettoyer, ne coche pas tout à l'aveugle dans WeChat, la galerie photo ou les documents ; pour les caches système, privilégie autant que possible l'outil de stockage intégré de la marque, les outils tiers ne servant qu'en complément.
  \item Nettoyage de stockage intégré Public : tout le monde, en particulier les seniors et ceux qui veulent éviter les complications. Points forts : Android [Réglages] > [Stockage] > [Nettoyage] ; iOS [Réglages] > [Général] > [Stockage iPhone]. Le système connaît mieux la structure de cache de ses applications ; risque de suppression erronée minimal — premier choix.
  \item Gestion du stockage intégrée WeChat/QQ Public : gros utilisateurs de messagerie. Points forts : WeChat [Moi] > [Réglages] > [Général] > [Espace de stockage], nettoyage fin par discussion et type de fichier ; plus sûr que les outils externes de « nettoyage WeChat ».
  \item Principe : utilisez d'abord ce qui est intégré au système ou à l'application ; si vous pouvez cocher par chemin/type, évitez le « tout effacer en un clic ». L'objectif est de libérer de l'espace, pas de balayer vos données importantes.
\end{itemize}

\booksubsection{Anti-indésirables et sécurité}

\begin{itemize}
  \item Centre national anti-fraude Public : tout le monde, surtout les seniors, les métiers liés à la finance et les acheteurs en ligne / destinataires d'appels d'inconnus. Points forts : solution officielle, alertes d'appels, SMS à risque, détection d'URL suspectes ; liée aux bases anti-fraude de la police — le maillon le plus fiable des « trois lignes de défense ».
  \item Centre de sécurité de la marque / gestionnaire de téléphone Public : utilisateurs de la marque concernée. Points forts : Huawei Phone Manager, Xiaomi Security Center, OPPO/vivo Phone Manager, etc. : blocage des appels indésirables, antivirus, gestion des autorisations, protection des paiements ; souvent gratuits et mieux intégrés au système. Évitez d'installer une seconde suite de sécurité commerciale qui créerait des conflits d'autorisations.
  \item Services anti-indésirables de l'opérateur Public : utilisateurs bombardés d'appels sortants ou vocaux robots. Points forts : China Mobile, envoyez KTFSR au 10086 ; China Telecom, suivez « Tianyi Anti-Spam » ; China Unicom, configurez dans « Unicom Mobile Manager ». Le blocage a lieu côté réseau : le téléphone ne sonne pas et une partie des numéros indésirables est déjà filtrée.
  \item AdGuard Public : tous les utilisateurs harcelés par pop-ups, bannières et publicités in-app. Points forts : blocage au niveau système (navigateur et au-delà), filtrage des bannières et du pistage au démarrage, économie de données et pages plus rapides ; version payante multi-appareils.
  \item Méfiez-vous des suites « nettoyage / anti-indésirables » : certaines tournent en arrière-plan, demandent sans cesse les contacts et les SMS, sans être plus efficaces que le système — elles consomment batterie et exposent votre vie privée. Activez d'abord les fonctions intégrées, puis envisagez l'application anti-fraude officielle.
\end{itemize}

\textbf{Principes de choix des logiciels pratiques}\par

\begin{itemize}
  \item Le système d'abord : blocage, nettoyage, autorisations — commencez par les interrupteurs intégrés.
  \item Complément officiel : Centre national anti-fraude et protection de l'opérateur peuvent se cumuler.
  \item Tiers avec retenue : une seule application par fonction, pour éviter les conflits d'autorisations et la concurrence en arrière-plan.
  \item Moins, c'est plus : plus vous installez d'outils, plus les notifications et autorisations se multiplient — et plus ces outils eux-mêmes risquent de vous déranger.
\end{itemize}

\booksubsection{Blocage intégré des appels indésirables : activation par marque}

Chemins d'accès par marque ci-dessous. Les libellés peuvent varier selon la version ; recherchez « indésirable », « blocage » ou « appel inconnu » dans les réglages. Complément du chapitre 8 « Gestion des appels indésirables » : ici, l'emplacement des interrupteurs ; là-bas, le signalement des numéros et les contestations.

\begin{center}
\begin{tabularx}{\linewidth}{|X|X|X|}
\hline
\textbf{Plateforme} & \textbf{Chemin d'accès} & \textbf{Options recommandées} \tabularnewline
\hline
iPhone / iOS & {}[Réglages] > [Téléphone] / [Messages] & Silence des appels inconnus ; contacts bloqués ; « Filtrer les expéditeurs inconnus » dans Messages \tabularnewline
Xiaomi / Redmi & Application Téléphone > Réglages > Blocage des appels indésirables ; ou Centre de sécurité & Blocage intelligent, coupe des appels d'un son, filtrage SMS par mots-clés \tabularnewline
Huawei / Honor & Téléphone > Blocage des appels indésirables ; Phone Manager > Blocage des appels indésirables & Identification des appels, liste noire, filtrage intelligent des SMS \tabularnewline
OPPO / realme / OnePlus & Application Téléphone > Blocage des appels indésirables ; ou Appareil > Blocage des appels indésirables & Blocage d'appels, identification des arnaques, règles de filtrage SMS \tabularnewline
vivo / iQOO & Téléphone > Blocage des appels indésirables ; i Manager & Blocage des appels indésirables, alertes d'arnaques, filtrage des messages \tabularnewline
Samsung & Téléphone > Réglages > Bloquer les numéros / Identification d'appel et appels indésirables & Bloquer les numéros inconnus, identification des appels indésirables (selon la région) \tabularnewline
\hline
\end{tabularx}
\end{center}

\begin{tcbinfo}{Superposer trois couches de protection}
Au quotidien, combinez « blocage système + protection opérateur + Centre national anti-fraude » : le système gère le signalement et la liste noire, l'opérateur bloque côté réseau les appels à haute fréquence, le centre anti-fraude apporte les alertes officielles. Inutile de tout régler au maximum : évitez de bloquer livraisons, repas à emporter ou appels médicaux — autorisez vos contacts et les personnes récemment contactées.

\end{tcbinfo}

\begin{tcbpractice}{Exercice pratique : activez le blocage dès aujourd'hui}
\begin{enumerate}
  \item Dans les réglages de l'application Téléphone, trouvez « Blocage des appels indésirables / Silence des appels inconnus » et activez le blocage intelligent.
  \item Faites de même sur le téléphone de vos parents : même chemin + service anti-indésirables de l'opérateur (China Mobile : envoyez KTFSR au 10086).
  \item Installez « Centre national anti-fraude » depuis la boutique d'applications, créez un compte et activez les alertes d'appels.
  \item Testez avec un numéro de test déjà signalé (ou consultez l'historique de blocage) pour vérifier que la fonction fonctionne.
\end{enumerate}

\end{tcbpractice}

\booksection{VII. Alertes météo, catastrophes naturelles et séismes}

Le blocage des appels protège des « arnaques » ; les alertes protègent des « intempéries et catastrophes ». Beaucoup n'installent qu'une application météo pour la température, sans notifications ni alerte sismique — or, dans les secondes qui précètent une alerte rouge de pluie torrentielle, un typhon ou un séisme, le téléphone peut vous faire gagner un temps précieux.

\booksubsection{Alertes météo : pluies torrentielles, typhons, fortes chaleurs}

\begin{itemize}
  \item Canaux officiels d'abord : les alertes de l'Administration météorologique de Chine, du Centre météorologique national et des observatoires provinciaux font autorité. Suivez le compte WeChat ou Weibo « China Weather », ou installez une application officielle comme « China Weather ».
  \item Applications Météo système : Huawei Weather, Xiaomi Weather, OPPO Weather, vivo Weather, Météo Apple, etc. poussent en général les alertes locales. Ouvrez leurs réglages et vérifiez que « Alertes météo / Intempéries » est activé.
  \item Caiyun Weather et prévisions à la minute : utiles avant un trajet ou une sortie d'école pour savoir s'il faut un parapluie dans la demi-heure ; les alertes officielles restent la référence.
  \item Couleurs d'alerte (usage général) : Bleu : restez attentif, déplacements normaux, suivez les mises à jour.
  \item Jaune : renforcez la prudence, limitez les activités en extérieur non nécessaires.
  \item Orange : grave, restez à l'intérieur si possible et éloignez-vous des zones dangereuses.
  \item Rouge : particulièrement grave, suivez les consignes de fermeture des écoles/chantiers et d'évacuation ; ne vous engagez pas dans les zones inondées, les torrents ou les vents violents.
\end{itemize}

Types d'alertes courants : pluies torrentielles, neige abondante, typhons, fortes chaleurs, vague de froid, vents violents, orages, grêle, verglas, brouillard, sécheresse, risques de catastrophes géologiques.

\booksubsection{Alertes de catastrophes naturelles : pas seulement la pluie}

\begin{itemize}
  \item Émetteurs : les observatoires pour la météo ; les risques géologiques et torrents sont souvent publiés conjointement par les ressources naturelles, l'eau et la protection civile ; les autorités locales diffusent aussi par radio d'urgence, SMS, 12379, etc.
  \item Recevoir les alertes sur plusieurs canaux : Notifications du système + alertes de l'application Météo.
  \item Suivre sur WeChat les comptes officiels locaux « Protection civile » et « Observatoire ».
  \item Canaux 12379 de diffusion des alertes nationales (SMS/service selon la région).
  \item Télévision, radio d'urgence, haut-parleurs de village/quartier — en montagne ou en saison des pluies, ne comptez pas uniquement sur le téléphone.
\end{itemize}

En extérieur et en déplacement : par typhon, évitez les panneaux et les arbres ; par fortes pluies, pas de garages souterrains, tunnels ou ponts bas ; en montagne, une alerte « risque météo de catastrophe géologique » impose de quitter immédiatement les pentes raides et les vallées.

Lien avec le manuel d'urgence : en cas de perte de contact, de batterie vide ou de besoin de géolocalisation après une catastrophe, voir le chapitre 32 (séismes/catastrophes) et le chapitre 7 (communication satellite et hors réseau).

\booksubsection{Alerte sismique : ce n'est pas une « prévision »}

L'alerte sismique (EEW) envoie, quelques secondes à quelques dizaines de secondes avant l'arrivée des ondes destructrices, une alerte aux zones potentiellement touchées, en s'appuyant sur la vitesse des ondes électromagnétiques et la détection précoce des ondes P. Elle ne prédit pas un séisme, mais offre une fenêtre pour « se protéger, couper le feu, sortir de l'ascenseur ». Plus on est proche de l'épicentre, moins le délai est long — parfois insuffisant.

\begin{itemize}
  \item Solutions officielles et sectorielles : en Chine continentale, on trouve le réseau national d'alerte sismique et les services portés par l'Institut de réduction des catastrophes de Chengdu ; certains téléphones intègrent ou reçoivent ces alertes en ligne.
  \item Applications tierces : cherchez « alerte sismique » dans la boutique, avec des applications officielles ou d'organismes reconnus (vérifiez l'éditeur avant l'installation ; méfiez-vous des fausses alertes qui demandent des autorisations abusives).
  \item Que faire à la réception de l'alerte : à l'intérieur, protégez-vous près d'une table, d'un angle porteur ou d'un meuble solide ; couvrez tête et cou ; éloignez-vous des fenêtres, murs extérieurs et meubles hauts ; ne fuyez pas aveuglément par les escaliers ni ne sautez ; quittez l'ascenseur à l'étage le plus proche ; dehors, gagnez un espace dégagé, loin des câbles, panneaux et bâtiments fragiles. Après le séisme, vérifiez gaz et électricité, attention aux répliques.
  \item Ne coupez pas les alertes : l'alerte sismique doit pouvoir « percer » le mode silencieux/Ne pas déranger ; la nuit, gardez le téléphone à portée de main et ne mettez pas le volume à zéro.
\end{itemize}

\booksubsection{Alerte sismique intégrée : activation par marque}

Les chemins peuvent varier selon la version ; cherchez « séisme », « alerte » ou « alerte sécurité » dans les réglages. iOS prend peu en charge les alertes sismiques en Chine : installez une application d'alerte fiable et autorisez localisation et notifications.

\begin{center}
\begin{tabularx}{\linewidth}{|X|X|X|}
\hline
\textbf{Plateforme} & \textbf{Chemin courant} & \textbf{Remarques} \tabularnewline
\hline
Xiaomi / Redmi & Phone Manager / Centre de sécurité > Alerte sismique ; ou recherche « séisme » dans Réglages & Certaines versions intègrent le service ; connexion et localisation requises \tabularnewline
Huawei / Honor & Réglages > Sécurité > Notifications d'alerte d'urgence / Alerte sismique & Peut s'appeler « Notifications d'alerte d'urgence » et inclure les séismes \tabularnewline
OPPO / realme / OnePlus & Réglages > Sécurité (SOS urgence) > Alerte sismique & Selon le modèle et la couverture régionale du service \tabularnewline
vivo / iQOO & Application Météo ou Réglages > Sécurité/Plus > Alerte sismique & Se fier aux menus de la version système réelle \tabularnewline
iPhone / iOS & Interrupteur intégré souvent incomplet ; installez une application d'alerte fiable et autorisez les notifications & L'application doit pouvoir tourner en arrière-plan et accéder à la localisation ; n'en bloquez pas les alertes \tabularnewline
Samsung et autres & Réglages > Sécurité et situations d'urgence / Alertes d'urgence sans fil & La couverture des « alertes d'urgence » varie selon la région et l'opérateur \tabularnewline
\hline
\end{tabularx}
\end{center}

\begin{tcbwarn}{Idées reçues sur l'alerte sismique}
\begin{itemize}
  \item Alerte $\neq$ prévision : elle ne dit pas « y aura-t-il un séisme demain », elle n'alerte qu'une fois le séisme déjà survenu et avant l'arrivée des ondes.
  \item Pas toujours plusieurs secondes d'avance : près de l'épicentre, il ne reste que quelques secondes, voire moins ; l'alerte sert surtout les zones périphériques.
  \item Couper les notifications = ne rien avoir : même la meilleure alerte échoue si le téléphone est en silencieux, si l'application est tuée ou si la localisation est refusée.
  \item Fausses applications : méfiez-vous des logiciels de « prévision sismique » qui demandent contacts et SMS ; privilégiez le système ou les organismes reconnus.
\end{itemize}

\end{tcbwarn}

\begin{tcbinfo}{Bien combiner les capacités d'alerte}
Au quotidien : « notifications Météo système + canaux météo/urgence officiels + alerte sismique (système ou application fiable) ». Comme pour le anti-indésirables, inutile d'empiler les logiciels : l'essentiel est d'autoriser les notifications et la localisation, et de laisser les alertes percer le mode Ne pas déranger. Après avoir configuré le blocage sur le téléphone de vos parents, vérifiez aussi météo et alerte sismique.

\end{tcbinfo}

\begin{tcbpractice}{Exercice pratique : 10 minutes pour vérifier les alertes}
\begin{enumerate}
  \item Ouvrez l'application Météo système, ses réglages, et activez les notifications « Alertes météo / Intempéries ».
  \item Dans les réglages système, cherchez « séisme » ou « alerte » et activez l'alerte sismique/d'urgence (Android) ; sur iPhone, installez une application d'alerte fiable et autorisez les notifications.
  \item Suivez les comptes WeChat/services de l'observatoire local ou de la protection civile, et autorisez leurs notifications.
  \item Vérifiez que le son d'alerte n'est pas totalement coupé par le mode Ne pas déranger ; faites de même sur le téléphone de vos parents.
\end{enumerate}

\end{tcbpractice}

\booksection{VIII. Collaboration multi-appareils : contrôle à distance, transfert rapide et sécurité des comptes}

Un téléphone vraiment pratique, c'est souvent un téléphone « en équipe » avec l'ordinateur et la tablette. Voici quelques outils très utiles : contrôle d'un ordinateur depuis le téléphone, bureautique IA et collaboration à distance, fonds d'écran animés, transfert local rapide, presse-papiers multi-appareils, codes dynamiques 2FA, numérisation papier et planification. Pour une vision plus systématique, voir le chapitre 17.

\booksubsection{Contrôler un ordinateur depuis le téléphone : Sunlogin / ToDesk}

\begin{itemize}
  \item Sunlogin Public : ceux qui doivent piloter un Windows/Mac à la maison ou au bureau depuis l'extérieur ; aussi utile pour dépanner à distance le téléphone des parents (voir chapitre 27). Points forts : écosystème chinois mature, bureau distant depuis le téléphone, transfert de fichiers, allumage/extinction (matériel parfois requis) ; la version gratuite suffit à un usage personnel léger.
  \item ToDesk Public : utilisateurs exigeant stabilité et faible latence pour le télétravail ou l'administration. Points forts : connexion par code d'appareil + mot de passe, téléphone vers ordinateur et inversement (selon version), transfert de fichiers et multi-écrans fluides ; nœuds souvent performants en Chine.
  \item Usage de base : installez le client sur l'ordinateur et connectez-vous / restez en ligne installez l'application sur le téléphone ouvrez la session avec l'identifiant d'appareil + le mot de passe (ou le même compte) pilotez le bureau avec le pavé tactile/gestes du téléphone.
  \item Règles de sécurité (à retenir) : Si un inconnu se faisant passer pour le « service client / la police / l'opérateur » vous demande d'installer Sunlogin ou ToDesk pour contrôler votre téléphone ou votre PC, refusez toujours — c'est une arnaque typique (voir chapitre 27).
  \item Pour votre propre usage : activez la double vérification, changez les mots de passe par défaut, déconnectez-vous ou désactivez l'accès sans surveillance quand vous n'utilisez pas la session ; sur un ordinateur d'entreprise, respectez la politique de sécurité interne.
  \item En télétravail dans un lieu public, vérifiez que personne ne peut lire votre écran ou vos codes.
\end{itemize}

\booksubsection{Bureautique IA : Trae / WorkBuddy / Dumate}

\begin{itemize}
  \item Trae — Pour qui : ceux qui collaborent sur du code et des documents, avec un environnement sinophone très accueillant. Points forts : un IDE IA / assistant de programmation signé ByteDance, dont le mode SOLO génère un projet à partir d'une seule phrase ; très efficace au quotidien pour écrire du code et retoucher des documents. Si tu l'utilises sur ordinateur, tu peux suivre et gérer les tâches depuis le mobile via une prise en main à distance ou les fonctions officielles associées (selon la forme actuelle du produit). Voir le Chapitre 41, Panorama des outils de programmation IA.
  \item WorkBuddy Public : professionnels voulant que l'IA rédige comptes rendus, e-mails, présentations et documents. Points forts : assistant bureautique IA pour les tâches « rédiger — polir — résumer » ; certains flux peuvent être lancés sur mobile, le gros du travail se faisant sur ordinateur.
  \item Dumate Public : équipes ou individus ayant besoin de collaboration et de gestion unifiée des documents et connaissances. Points forts : outil collaboratif IA centré sur le rassemblement des contenus et l'édition partagée ; consultation et retouches légères sur mobile, mise en page complexe sur grand écran.
  \item Comprendre « contrôler un PC depuis le téléphone » : de nombreuses fonctions bureautiques IA tournent sur ordinateur ou le Web. Depuis le téléphone, via Sunlogin/ToDesk ou l'application officielle, vous « télécommandez » les tâches IA du PC pendant les trajets. Principe : calculs lourds et longs documents sur l'ordinateur ; suivi et petites retouches sur le téléphone.
  \item Vie privée : avant d'envoyer contrats, pièces d'identité ou données clients, lisez la politique de confidentialité et les règles de conformité ; pour les documents d'entreprise, utilisez d'abord les outils imposés par l'organisation.
\end{itemize}

\booksubsection{Bureaux animés : fonds d'écran miHoYo (mobile + PC)}

\begin{itemize}
  \item Rengong Zhuomian — Pour qui : les fans des personnages de miHoYo et celles et ceux qui veulent des fonds d'écran animés gratuits. Points forts : une appli de fonds d'écran animés gratuite signée miHoYo, disponible sur mobile et PC ; les personnages en Live2D ou en animation peuvent habiller ton bureau, et sur certains modèles ou systèmes être définis comme écran de verrouillage ou fond d'écran.
  \item Autonomie et performance : un fond animé consomme plus de batterie et de GPU ; en cas de batterie faible ou d'ancien modèle, utilisez-le uniquement pendant la charge ou réduisez la fréquence d'images si l'application le permet.
  \item Source d'installation : privilégiez le site officiel ou les boutiques reconnues ; évitez les « versions craquées » qui contiennent souvent des chevaux de Troie ; sur PC, vérifiez le chemin d'installation et le lancement au démarrage.
\end{itemize}

\booksubsection{Transfert de fichiers en réseau local : LocalSend / Feige / AirDrop}

\begin{itemize}
  \item LocalSend Public : mélanges Android/Windows/iOS/Mac/Linux, transfert multi-plateformes sans cloud public. Points forts : open source, connexion directe en réseau local, partage de photos, vidéos et documents sur le même Wi-Fi, sans câble ni serveur constructeur ; vitesse selon le routeur.
  \item Feige Chuanshu — Pour qui : les échanges entre PC Windows de bureau et les environnements de travail en réseau local un peu anciens. Points forts : un outil vétéran de transfert et de messagerie en réseau local, idéal pour faire circuler de gros fichiers en intranet ; télécharge-le depuis un canal officiel et respecte la politique informatique de ton entreprise sur son réseau interne.
  \item AirDrop Public : utilisateurs iPhone, iPad, Mac. Points forts : partage en un clic dans l'écosystème Apple, découverte Bluetooth + transfert Wi-Fi direct, sans consommation de données ; mode « Contacts uniquement » pour limiter les sollicitations.
  \item Autres solutions : Mi Share, Huawei Share, vivo Share, etc. sont très rapides entre appareils de même marque ; l'assistant de transfert de fichiers WeChat est pratique mais limité en taille et en appareils connectés — pour les gros fichiers, préférez LocalSend ou AirDrop.
  \item Conseils : les deux appareils doivent être sur le même Wi-Fi ; évitez le Wi-Fi public pour les fichiers sensibles, préférez un point d'accès privé ; après l'envoi, vérifiez l'intégrité du fichier côté récepteur.
\end{itemize}

\booksubsection{Presse-papiers multi-appareils : KDE Connect / transfert téléphone-ordinateur}

\begin{itemize}
  \item Presse-papiers universel Apple Public : utilisateurs iPhone + Mac. Points forts : avec le même Apple ID et Handoff activé, copiez sur le téléphone collez sur le Mac (et inversement), texte et certains fichiers.
  \item KDE Connect Public : Android + Windows/Linux/Mac, synchronisation open source, miroir de notifications, saisie à distance. Points forts : après appairage local, synchronisation du presse-papiers, notifications du téléphone sur l'ordinateur, téléphone comme télécommande/touchpad de présentation ; les données ne passent pas par un compte tiers.
  \item Solutions de marque « téléphone-ordinateur » : Huawei Multi-screen Collaboration / Super Device, Xiaomi Miaoxiang, vivo Share, OPPO Connect, Honor Connect, etc. offrent souvent presse-papiers inter-écrans et glisser-déposer ; l'expérience est optimale au sein d'une même marque.
  \item Canal officiel Windows + Android : « Lien vers Windows / Link to Windows » (évolution de Your Phone) affiche notifications et photos du téléphone sur le PC, avec copier-coller dans les cas pris en charge ; intégration plus poussée sur Samsung, etc.
  \item Sécurité : le presse-papiers peut contenir mots de passe et codes ; évitez la synchronisation globale sur un ordinateur partagé et ne recollez pas dans les champs de mot de passe ; certaines claviers lisent aussi le presse-papiers — surveillez les autorisations.
\end{itemize}

\booksubsection{Codes dynamiques 2FA : Aegis / Microsoft Authenticator}

\begin{itemize}
  \item Pourquoi un authentificateur plutôt que le SMS : un code SMS peut être détourné par un remplacement de carte SIM ; un code TOTP est généré localement selon l'heure, sans dépendre du réseau mobile, plus résistant à l'hameçonnage et à l'interception (mais pas aux malwares ni aux captures d'écran).
  \item Aegis Authenticator Public : utilisateurs Android préférant l'open source et le contrôle local. Points forts : gratuit et open source, sauvegardes chiffrées, import/export, interface propre, déverrouillage par empreinte/visage de la liste des jetons.
  \item Microsoft Authenticator Public : utilisateurs de comptes Microsoft, Office 365 ou d'entreprise, sur iOS/Android. Points forts : validation par notification push, surveillance des mots de passe et sauvegarde de comptes (selon version), forte intégration avec Windows et les identités d'entreprise.
  \item Google Authenticator / 2FAS, etc. : également courants dans ces écosystèmes ; l'essentiel est d'avoir une sauvegarde chiffrée pour ne perdre aucun jeton lors d'un changement de téléphone.
  \item Changement de téléphone : migrez depuis l'ancien appareil ou utilisez les « codes de secours/clé de récupération » ; imprimez les codes de secours et conservez-les à deux endroits, pas seulement sur le même téléphone.
  \item Avec un gestionnaire de mots de passe : certains coffres stockent aussi les TOTP — pratique, mais le risque se concentre ; pour les comptes précieux, préférez « gestionnaire de mots de passe + authentificateur distinct ». Voir le chapitre 14.
\end{itemize}

\booksubsection{Combinaison recommandée : une journée type}

\textbf{Exemple de flux de travail}\par

\begin{itemize}
  \item Matin : TickTick pour l'emploi du temps jeter un œil à la météo/alertes ouvrir Trae/documents sur l'ordinateur, téléphone posé.
  \item Journée : numériser les documents papier en PDF, les envoyer via LocalSend/AirDrop sur l'ordinateur ; rédiger les longs textes avec les outils IA sur PC.
  \item En déplacement : ToDesk/Sunlogin pour reprendre la main sur le PC ; connexions aux comptes importants via codes dynamiques Aegis/Microsoft Authenticator.
  \item Soir : cocher les tâches terminées dans TickTick ; changer de fond d'écran animé pour se détendre si besoin.
\end{itemize}

\begin{tcbwarn}{Règles de sécurité pour le contrôle à distance et les comptes}
\begin{itemize}
  \item Un inconnu vous demande d'installer un outil de contrôle à distance raccrochez et contactez vous-même le service client officiel.
  \item Ne donnez de code de vérification ou de code dynamique à personne, même à un « livreur » ou un « service client ».
  \item Avant de contrôler un PC d'entreprise à distance ou d'y installer des logiciels IA, demandez l'accord de l'IT.
  \item Restez très vigilant lors de transferts de fichiers ou de connexions à distance sur un réseau public.
\end{itemize}

\end{tcbwarn}

\begin{tcbpractice}{Exercice pratique : faire fonctionner une fois la chaîne multi-appareils}
\begin{enumerate}
  \item Connectez téléphone et ordinateur au même Wi-Fi et envoyez-vous une photo avec LocalSend (ou AirDrop).
  \item Créez dans TickTick la tâche « Ce soir 20 h : trier les scans » et vérifiez qu'elle apparaît sur le téléphone et l'ordinateur.
  \item Activez la 2FA sur votre e-mail ou WeChat (si disponible), liez Aegis ou Microsoft Authenticator, et conservez les codes de secours.
  \item Installez Sunlogin ou ToDesk sur votre propre ordinateur, connectez-vous une fois depuis votre téléphone, puis déconnectez-vous ou renforcez le mot de passe — maîtrisez le processus sans laisser de porte dérobée.
\end{enumerate}

\end{tcbpractice}
